\subsubsection{Data Management}
\begin{enumerate}
    \item \textbf{Data Reading and Processing} \\
          \begin{enumerate}
            \item Reading Data
                  \begin{itemize}
                    \item \textbf{Read Gene IDs from TSV File}
                          
                          Reads gene identifiers from a TSV file containing expression data.
                          Typically reads from the first column after skipping header rows. Call
                          the \texttt{read\_gene\_ids\_from\_tsv\_file} subroutine with the
                          following arguments:
                          
                          \textbf{Input arguments:}
                          \begin{itemize}
                            \item \texttt{filename character(len=*)}: Path to the TSV file
                                  
                            \item \texttt{n\_header\_rows integer(int32)}: Number of header rows
                                  to skip (typically 1)
                                  
                            \item \texttt{gene\_col integer(int32)}: Column index containing
                                  gene IDs (typically 1)
                          \end{itemize}
                          \textbf{Output arguments:}
                          \begin{itemize}
                            \item \texttt{gene\_ids character(len=*)}: Pre-allocated array for
                                  gene identifiers
                                  
                            \item \texttt{ierr integer(int32)}: Error code
                          \end{itemize}
                          \begin{lstlisting}[language=Fortran, caption={Read Gene IDs from TSV File}]
character(len=256) :: gene_ids(n_genes)  
! Pre-allocate with expected size

integer(int32) :: ierr

call read_gene_ids_from_tsv_file('expression_data.tsv', &
                             gene_ids, 1, 1, ierr)
                          \end{lstlisting}
                          
                    \item \textbf{Read Expression Vectors}
                          
                          Reads expression data from one or more TSV/CSV files and populates
                          a pre-allocated 2D array (samples x genes). Uses hashmap for
                          efficient gene ID matching and validates for finite values. Call the
                          \texttt{read\_expression\_vectors} subroutine with the following
                          arguments:
                          
                          \textbf{Input arguments:}
                          \begin{itemize}
                            \item \texttt{file\_list character(len=*)}: Array of file paths to
                                  read from
                                  
                            \item \texttt{gene\_ids character(len=*)}: Reference gene IDs for
                                  validation and matching
                                  
                            \item \texttt{expression\_vectors real(real64)}: Pre-allocated 2D
                                  array for expression data
                                  
                            \item \texttt{n\_header\_rows integer(int32)}: Number of header rows
                                  to skip
                                  
                            \item \texttt{gene\_col integer(int32)}: Column index containing
                                  gene IDs
                                  
                            \item \texttt{value\_cols integer(int32)}: Array of column indices
                                  containing expression values
                                  
                            \item \texttt{start\_row integer(int32)}: Starting row index in expression\_vectors
                                  array
                                  
                            \item \texttt{delimiter character(len=1), optional}: Column delimiter
                                  (default: tab)
                          \end{itemize}
                          \textbf{Output arguments:}
                          \begin{itemize}
                            \item \texttt{expression\_vectors real(real64)}: Pre-allocated 2D
                                  array for expression data populated with data
                                  
                            \item \texttt{ierr integer(int32)}: Error code
                          \end{itemize}
                          \begin{lstlisting}[language=Fortran, caption={Read Expression Vectors}]
! Pre-allocate arrays
real(real64) :: expr_data(n_samples, n_genes)

! Read TSV file with tab delimiter (default)
call read_expression_vectors(files, gene_ids, expr_data, 1, 1, &
                            value_cols, 1, ierr)

! Read CSV file with comma delimiter
call read_expression_vectors(files, gene_ids, expr_data, 1, 1, &
                            value_cols, 1, ierr, ',')
                          \end{lstlisting}
                          
                    \item \textbf{Read OrthoFinder Family Mapping}
                          
                          Reads gene family assignments from OrthoFinder output TSV file and
                          creates mapping arrays between genes and families. The \texttt{gene\_to\_fam}
                          array contains for each gene the index of it's family. Uses
                          hashmap for efficient gene lookup. Call the \texttt{read\_orthofinder\_file}
                          subroutine with the following arguments:
                          
                          \textbf{Input arguments:}
                          \begin{itemize}
                            \item \texttt{filename character(len=*)}: Path to OrthoFinder TSV
                                  file
                                  
                            \item \texttt{gene\_ids character(len=*)}: Pre-allocated array of
                                  gene IDs to match
                          \end{itemize}
                          \textbf{Output arguments:}
                          \begin{itemize}
                            \item \texttt{family\_ids character(len=*)}: Pre-allocated array
                                  for family identifiers
                                  
                            \item \texttt{gene\_to\_fam integer(int32)}: Pre-allocated array
                                  for gene to family mapping
                                  
                            \item \texttt{ierr integer(int32)}: Error code
                          \end{itemize}
                          \begin{lstlisting}[language=Fortran, caption={Read OrthoFinder Family Mapping}]
character(len=256) :: family_ids(n_families)  
! Pre-allocate with expected family count

integer(int32) :: gene_to_fam(n_genes)     
! Pre-allocate with gene count

integer(int32) :: ierr

call read_orthofinder_file('Orthogroups.tsv', gene_ids, &
                            family_ids, gene_to_fam, ierr)
                          \end{lstlisting}
                          
                          \textbf{Important Implementation Details:}
                          \begin{itemize}
                            \item All output arrays must be pre-allocated with the correct dimensions
                                  and sizes
                                  
                            \item Uses hashmap data structure for O(1) gene ID lookups
                                  
                            \item Validates for finite values (rejects NaN and Inf) in expression
                                  data
                                  
                            \item Handles both TSV (tab-separated) and CSV (comma-separated)
                                  formats, gene\_to\_fam array uses 1-based indexing for family
                                  assignments (0 indicates unassigned)
                          \end{itemize}
                          
                          \textbf{Typical Data Reading Workflow:}
                          \begin{lstlisting}[language=Fortran, caption={Complete Data Reading Workflow}]
! 1. Pre-allocate arrays

! 2. Read gene IDs from expression file
call read_gene_ids_from_tsv_file('expression.tsv', gene_ids, &
                                    1, 1, ierr)

! 3. Read expression data (columns 2-68 contain sample data)
character(len=256) :: files(1) = ['expression.tsv']

call read_expression_vectors(files, gene_ids, expr_data, 1, 1, &
                             samples, 1, ierr)

! 4. Read family mapping from OrthoFinder
call read_orthofinder_file('Orthogroups.tsv', gene_ids, & 
                            family_ids, gene_to_fam, ierr)
                          \end{lstlisting}
                  \end{itemize}
          \end{enumerate}
          
    \item \textbf{Data Filtering and Validation}
          \begin{itemize}
            \item \textbf{Get Unassigned Genes Mask}
                  
                  Creates a logical mask identifying genes that have no family assignment
                  (gene\_to\_fam value $=$ 0). Call the \texttt{get\_unassigned\_mask} subroutine
                  with the following arguments:
                  
                  \textbf{Input arguments:}
                  \begin{itemize}
                    \item \texttt{gene\_to\_fam integer(int32)}: Gene to family mapping array
                  \end{itemize}
                  \textbf{Output arguments:}
                  \begin{itemize}
                    \item \texttt{mask logical}: Logical mask array (same size as gene\_to\_fam)
                          
                    \item \texttt{n\_genes\_kept integer(int32)}: Number of genes with family
                          assignments
                          
                    \item \texttt{ierr integer(int32)}: Error code
                  \end{itemize}
                  \begin{lstlisting}[language=Fortran, caption={Create Unassigned Genes Mask}]
logical :: unassigned_mask(n_genes)
integer(int32) :: n_kept, ierr

call get_unassigned_mask(gene_to_fam, unassigned_mask, n_kept)
! unassigned_mask(i) = .true. if gene i has family assignment
! n_kept = count of genes with family assignments
                  \end{lstlisting}
                  
            \item \textbf{Apply Unassigned Mask}
                  
                  Filters out genes without family assignments from all data arrays (gene\_ids,
                  expression\_vectors, gene\_to\_fam). Arrays are reallocated to the
                  filtered size. Call the \texttt{apply\_unassigned\_mask} subroutine
                  with the following arguments:
                  
                  \textbf{Input arguments:}
                  \begin{itemize}
                    \item \texttt{mask logical}: Logical mask from get\_unassigned\_mask
                          
                    \item \texttt{n\_genes\_kept integer(int32)}: Number of genes to keep
                          (from get\_unassigned\_mask)
                  \end{itemize}
                  \textbf{Input/Output arguments:}
                  \begin{itemize}
                    \item \texttt{gene\_ids character(len=*), allocatable}: Gene IDs array
                          (reallocated)
                          
                    \item \texttt{expression\_vectors real(real64), allocatable}: Expression
                          data matrix (reallocated)
                          
                    \item \texttt{gene\_to\_fam integer(int32), allocatable}: Gene to family
                          mapping (reallocated)
                          
                    \item \texttt{ierr integer(int32)}: Error code
                  \end{itemize}
                  \begin{lstlisting}[language=Fortran, caption={Filter Unassigned Genes}]
character(len=256), allocatable :: gene_ids(n_genes)
real(real64), allocatable :: expr_data(n_samples,n_genes)
integer(int32), allocatable :: gene_to_fam(n_genes)
logical :: unassigned_mask(n_genes)
integer(int32) :: n_kept, ierr

! First create the mask
call get_unassigned_mask(gene_to_fam, unassigned_mask, n_kept)

! Then apply it to filter all arrays
call apply_unassigned_mask(gene_ids, expr_data, gene_to_fam, &
                          unassigned_mask, n_kept, ierr)
! Arrays are now reallocated to size n_kept
                  \end{lstlisting}
                  
            \item \textbf{Validate Data Structure}
                  
                  Performs comprehensive validation of data dimensions and consistency across
                  all arrays. Call the \texttt{validate\_data\_structure} subroutine with
                  the following arguments:
                  
                  \textbf{Input arguments:}
                  \begin{itemize}
                    \item \texttt{n\_genes integer(int32)}: Number of genes
                          
                    \item \texttt{n\_families integer(int32)}: Number of gene families
                          
                    \item \texttt{n\_samples integer(int32)}: Number of samples
                          
                    \item \texttt{gene\_ids character(len=*)}: Gene identifier array
                          
                    \item \texttt{family\_ids character(len=*)}: Family identifier array
                          
                    \item \texttt{gene\_to\_family integer(int32)}: Gene to family mapping
                          
                    \item \texttt{expression real(real64)}: Expression data matrix
                          
                    \item \texttt{centroids real(real64)}: Family centroids matrix
                          
                    \item \texttt{shift\_vectors real(real64)}: Shift vectors matrix
                          
                    \item \texttt{ierr integer(int32)}: Error code
                  \end{itemize}
                  \begin{lstlisting}[language=Fortran, caption={Validate Data Structure}]
call validate_data_structure(n_genes, n_families, n_samples, &
                            gene_ids, family_ids, gene_to_fam, &
                            expr_data, centroids, shift_vectors, &
                            ierr)
                  \end{lstlisting}
                  
            \item \textbf{Validate Expression Data}
                  
                  Validates expression data matrix for NaN, Inf, and optionally negative
                  values. Call the \texttt{validate\_expression\_data} subroutine with the
                  following arguments:
                  
                  \textbf{Input arguments:}
                  \begin{itemize}
                    \item \texttt{expression real(real64)}: Expression data matrix to validate
                          
                    \item \texttt{allow\_negative logical}: Whether to allow negative expression
                          values
                          
                    \item \texttt{ierr integer(int32)}: Error code
                  \end{itemize}
                  \begin{lstlisting}[language=Fortran, caption={Validate Expression Data}]
! Validate expression data (disallow negative values)
call validate_expression_data(expr_data, .false., ierr)

! Validate allowing negative values (e.g., for log-fold changes)
call validate_expression_data(expr_data, .true., ierr)
                  \end{lstlisting}
                  
            \item \textbf{Validate Empty Strings}
                  
                  Validates that a string array contains no empty or whitespace-only entries.
                  Call the \texttt{validate\_empty\_strings} subroutine with the following
                  arguments:
                  
                  \textbf{Input arguments:}
                  \begin{itemize}
                    \item \texttt{string\_array character(len=*)}: String array to validate
                          
                    \item \texttt{array\_name character(len=*)}: Name for error messages
                  \end{itemize}
                  \textbf{Output arguments:}
                  \begin{itemize}
                    \item \texttt{ierr integer(int32)}: Error code
                  \end{itemize}
                  \begin{lstlisting}[language=Fortran, caption={Validate String Array}]
call validate_empty_strings(gene_ids, "Gene IDs", ierr)
call validate_empty_strings(family_ids, "Family IDs", ierr)
                  \end{lstlisting}
                  
            \item \textbf{Validate Gene to Family Mapping}
                  
                  Validates that all gene to family indices are within valid range (1 to
                  n\_families). Call the \texttt{validate\_gene\_to\_family\_mapping} subroutine
                  with the following arguments:
                  
                  \textbf{Input arguments:}
                  \begin{itemize}
                    \item \texttt{gene\_to\_family integer(int32)}: Gene to family mapping
                          array
                          
                    \item \texttt{n\_families integer(int32)}: Number of valid families
                  \end{itemize}
                  \textbf{Output arguments:}
                  \begin{itemize}
                    \item \texttt{ierr integer(int32)}: Error code
                  \end{itemize}
                  \begin{lstlisting}[language=Fortran, caption={Validate Gene-Family Mapping}]
call validate_gene_to_family_mapping(gene_to_fam, n_families, ierr)
                  \end{lstlisting}
                  
            \item \textbf{Validate String Array Uniqueness}
                  
                  Validates that all entries in a string array are unique (no duplicates).
                  Uses a hashet implementation internally for O(n) runtime. Call the
                  \texttt{validate\_string\_array\_uniqueness} subroutine with the following
                  arguments:
                  
                  \textbf{Input arguments:}
                  \begin{itemize}
                    \item \texttt{string\_array character(len=*)}: String array to check
                          for duplicates
                  \end{itemize}
                  \textbf{Output arguments:}
                  \begin{itemize}
                    \item \texttt{ierr integer(int32)}: Error code
                  \end{itemize}
                  \begin{lstlisting}[language=Fortran, caption={Validate Unique Strings}]
call validate_string_array_uniqueness(gene_ids, ierr)
call validate_string_array_uniqueness(family_ids, ierr)
                  \end{lstlisting}
          \end{itemize}
          
    \item \textbf{Data Accessors}
          \begin{itemize}
            \item \textbf{Get Gene Index}
                  
                  Finds the array index for a given gene identifier. Call the \texttt{get\_gene\_index}
                  function with the following arguments:
                  
                  \textbf{Input arguments:}
                  \begin{itemize}
                    \item \texttt{gene\_ids character(len=*)}: Array of gene identifiers
                          
                    \item \texttt{gene\_id character(len=*)}: Gene ID to find
                  \end{itemize}
                  \textbf{Return value:}
                  \begin{itemize}
                    \item \texttt{index integer(int32)}: Array index (1-based) or 0 if not
                          found
                  \end{itemize}
                  \begin{lstlisting}[language=Fortran, caption={Find Gene Index}]
integer(int32) :: gene_idx

gene_idx = get_gene_index(gene_ids, 'NP_001000001.1')
if (gene_idx > 0) then
    ! Gene found at position gene_idx
else
    ! Gene not found
end if
                  \end{lstlisting}
                  
            \item \textbf{Get Family Index}
                  
                  Finds the array index for a given family identifier. Call the \texttt{get\_family\_index}
                  function with the following arguments:
                  
                  \textbf{Input arguments:}
                  \begin{itemize}
                    \item \texttt{family\_ids character(len=*)}: Array of family identifiers
                          
                    \item \texttt{family\_id character(len=*)}: Family ID to find
                  \end{itemize}
                  \textbf{Return value:}
                  \begin{itemize}
                    \item \texttt{index integer(int32)}: Array index (1-based) or 0 if not
                          found
                  \end{itemize}
                  \begin{lstlisting}[language=Fortran, caption={Find Family Index}]
integer(int32) :: fam_idx

fam_idx = get_family_index(family_ids, 'OG0000001')
                  \end{lstlisting}
                  
            \item \textbf{Get Family for Gene Index}
                  
                  Retrieves the family index for a given gene index. Call the \texttt{get\_family\_for\_gene\_index}
                  subroutine with the following arguments:
                  
                  \textbf{Input arguments:}
                  \begin{itemize}
                    \item \texttt{gene\_index integer(int32)}: Gene array index (1-based)
                          
                    \item \texttt{gene\_to\_family integer(int32)}: Gene to family mapping
                          array
                  \end{itemize}
                  \textbf{Output arguments:}
                  \begin{itemize}
                    \item \texttt{family\_index integer(int32)}: Family array index (1-based)
                          
                    \item \texttt{ierr integer(int32)}: Error code
                  \end{itemize}
                  \begin{lstlisting}[language=Fortran, caption={Get Family for Gene}]
integer(int32) :: gene_idx, fam_idx, ierr

gene_idx = get_gene_index(gene_ids, 'NP_001000001.1')
call get_family_for_gene_index(gene_idx, gene_to_fam, fam_idx, ierr)
                  \end{lstlisting}
                  
            \item \textbf{Get Family Centroid}
                  
                  Retrieves the centroid expression profile for a given family. Call the
                  \texttt{get\_family\_centroid} subroutine with the following arguments:
                  
                  \textbf{Input arguments:}
                  \begin{itemize}
                    \item \texttt{family\_index integer(int32)}: Family array index (1-based)
                          
                    \item \texttt{family\_centroids real(real64)}: Family centroids matrix
                  \end{itemize}
                  \textbf{Output arguments:}
                  \begin{itemize}
                    \item \texttt{ierr integer(int32)}: Error code
                          
                    \item \texttt{centroid\_vector real(real64)}: Centroid expression profile
                  \end{itemize}
                  \begin{lstlisting}[language=Fortran, caption={Get Family Centroid}]
real(real64) :: centroid_vector(n_samples)
integer(int32) :: fam_idx, ierr

fam_idx = get_family_index(family_ids, 'OG0000001')
call get_family_centroid(fam_idx, centroids, centroid_vector, ierr)
                  \end{lstlisting}
                  
                  \textbf{Complete Validation Workflow:} 
                  \begin{lstlisting}[language=Fortran, caption={Complete Data Validation Workflow}]
integer(int32) :: ierr

! 1. Filter unassigned genes
logical :: unassigned_mask(n_genes)
integer(int32) :: n_kept
call get_unassigned_mask(gene_to_fam, unassigned_mask, n_kept)
call apply_unassigned_mask(gene_ids, expr_data, gene_to_fam, &
unassigned_mask, n_kept, ierr)

! 2. Validate individual components
call validate_empty_strings(gene_ids, "Gene IDs", ierr)
call validate_empty_strings(family_ids, "Family IDs", ierr)
call validate_string_array_uniqueness(gene_ids, ierr)
call validate_string_array_uniqueness(family_ids, ierr)
call validate_gene_to_family_mapping(gene_to_fam, n_families, ierr)
call validate_expression_data(expr_data, .false., ierr)

! 3. Comprehensive structure validation
call validate_data_structure(n_genes, n_families, n_samples, gene_ids, &
family_ids, gene_to_fam, expr_data, &
centroids, shift_vectors, ierr)

! 4. Use accessors for data retrieval
integer(int32) :: gene_idx, fam_idx
real(real64) :: centroid(n_samples)

gene_idx = get_gene_index(gene_ids, 'NP_001000001.1')
call get_family_for_gene_index(gene_idx, gene_to_fam, fam_idx, ierr)
call get_family_centroid(fam_idx, centroids, centroid, ierr)
            \end{lstlisting}
          \end{itemize}
\end{enumerate}